\documentclass[12pt,a4paper]{article}

% ---- Packages ----
\usepackage[margin=1in]{geometry}
\usepackage{graphicx}
\usepackage{listings}
\usepackage{xcolor}
\usepackage{hyperref}
\usepackage{amsmath}
\usepackage{booktabs}
\usepackage{fancyhdr}
\usepackage{titlesec}
\usepackage{float}
\usepackage{caption}

% ---- Colors ----
\definecolor{codegreen}{rgb}{0,0.6,0}
\definecolor{codepurple}{rgb}{0.58,0,0.82}
\definecolor{codegray}{rgb}{0.5,0.5,0.5}
\definecolor{backcolour}{rgb}{0.95,0.95,0.95}
\definecolor{darkblue}{rgb}{0.0,0.0,0.6}

% ---- Code Listings Style ----
\lstdefinestyle{verilogstyle}{
    backgroundcolor=\color{backcolour},
    commentstyle=\color{codegreen},
    keywordstyle=\color{darkblue}\bfseries,
    numberstyle=\tiny\color{codegray},
    stringstyle=\color{codepurple},
    basicstyle=\ttfamily\scriptsize,
    breaklines=true,
    numbers=left,
    numbersep=5pt,
    frame=single,
    tabsize=4,
    captionpos=b,
    morekeywords={module, endmodule, input, output, wire, reg, always, posedge, negedge, assign, begin, end, if, else, parameter, localparam, integer, for}
}
\lstset{style=verilogstyle}

% ---- Header / Footer ----
\pagestyle{fancy}
\fancyhf{}
\rhead{Assignment 11 -- Reset Synchronizer}
\lhead{Digital IC Design}
\rfoot{Page \thepage}

% ---- Title Format ----
\titleformat{\section}{\Large\bfseries\color{darkblue}}{\thesection}{1em}{}
\titleformat{\subsection}{\large\bfseries\color{darkblue!80}}{\thesubsection}{1em}{}
\titleformat{\subsubsection}{\normalsize\bfseries\color{darkblue!60}}{\thesubsubsection}{1em}{}

% ---- Hyperref ----
\hypersetup{
    colorlinks=true,
    linkcolor=darkblue,
    urlcolor=blue,
    citecolor=darkblue
}

\begin{document}

% ===========================================================================
% TITLE PAGE
% ===========================================================================
\begin{titlepage}
\centering
\vspace*{2cm}
{\Huge\bfseries Reset Synchronizer\\(RST\_SYNC)\par}
\vspace{1cm}
{\Large Assignment 11 -- Digital IC Design Course\par}
\vspace{2cm}
{\large\textbf{Submitted by:} Ziad Mostafa Abdelaziz\par}
\vspace{0.5cm}
{\large\today\par}
\vfill
\end{titlepage}

\tableofcontents
\newpage

% ===========================================================================
% 1. INTRODUCTION
% ===========================================================================
\section{Introduction}

In synchronous digital IC design, un-synchronized active-low asynchronous resets present severe timing hazards during reset release (de-assertion). If an asynchronous reset transitions from active ($0$) to inactive ($1$) near a clock edge, it can violate the \textbf{recovery time} ($t_{\text{recovery}}$) or \textbf{removal time} ($t_{\text{removal}}$) requirements of destination flip-flops, causing metastability and uncoordinated state transitions.

The \textbf{Reset Synchronizer (\texttt{RST\_SYNC})} solves this timing hazard by providing:
\begin{enumerate}
    \item \textbf{Immediate Asynchronous Reset Assertion:} Resets all destination flip-flops immediately when \texttt{RST} drops to $0$ without waiting for a clock edge.
    \item \textbf{Synchronous Reset De-assertion:} Releases the reset (\texttt{SYNC\_RST} $\to 1$) strictly on a rising edge of \texttt{CLK} after \texttt{NUM\_STAGES} clock cycles, preventing recovery/removal timing violations.
\end{enumerate}

% ===========================================================================
% 2. BLOCK INTERFACE
% ===========================================================================
\section{Block Interface}

\subsection{Parameters}
\begin{table}[H]
\centering
\begin{tabular}{@{}llc@{}}
\toprule
\textbf{Parameter Name} & \textbf{Description} & \textbf{Default Value} \\
\midrule
\texttt{NUM\_STAGES} & Number of D Flip-Flop Synchronizer Stages & 2 \\
\bottomrule
\end{tabular}
\caption{Design Parameters}
\end{table}

\subsection{Ports Description}
\begin{table}[H]
\centering
\begin{tabular}{@{}llc@{}}
\toprule
\textbf{Signal Name} & \textbf{Description} & \textbf{Width} \\
\midrule
\texttt{RST}      & Asynchronous active-low reset input & 1 \\
\texttt{CLK}      & Destination domain clock            & 1 \\
\texttt{SYNC\_RST} & Synchronized active-low reset output & 1 \\
\bottomrule
\end{tabular}
\caption{Port Descriptions}
\end{table}

% ===========================================================================
% 3. ARCHITECTURE AND CIRCUIT OPERATION
% ===========================================================================
\section{Architecture and Circuit Operation}

The \texttt{RST\_SYNC} module consists of a parameterized chain of \texttt{NUM\_STAGES} D flip-flops with active-low asynchronous reset inputs connected to the external \texttt{RST} line:

\begin{figure}[H]
\centering
\begin{verbatim}
       1 (VCC / Logic 1'b1)
       |
       v
     [D Q]_1 ---> [D Q]_2 ---> ... ---> [D Q]_NUM_STAGES ---> SYNC_RST
       ^            ^                     ^
      CLK          CLK                   CLK
       |            |                     |
     (Async)      (Async)               (Async)
       |            |                     |
       +------------+---------------------+ <--- RST (Active-Low)
\end{verbatim}
\caption{Reset Synchronizer Shift Register Architecture}
\end{figure}

\subsection{Circuit Dynamics}
\begin{itemize}
    \item \textbf{Assertion Phase ($\texttt{RST} = 0$):} All flip-flops in the chain are asynchronously reset to $0$. \texttt{SYNC\_RST} drops to $0$ immediately ($< 1\text{ ns}$ delay).
    \item \textbf{De-assertion Phase ($\texttt{RST} = 1$):} The asynchronous reset is released. Constant high logic ($1\text{'b1}$) shifts through the chain on each \texttt{posedge CLK}. After \texttt{NUM\_STAGES} clock edges, \texttt{SYNC\_RST} transitions to $1$ synchronously with \texttt{CLK}.
\end{itemize}

% ===========================================================================
% 4. RTL VERILOG CODE
% ===========================================================================
\section{RTL Verilog Module}

\lstinputlisting[language=Verilog, caption={RST\_SYNC.v -- Synthesizable Reset Synchronizer RTL}]{../rtl/RST_SYNC.v}

% ===========================================================================
% 5. TESTBENCH AND SIMULATION RESULTS
% ===========================================================================
\section{Testbench and Simulation Results}

\subsection{Testbench Implementation}
The self-checking testbench (\texttt{RST\_SYNC\_tb.v}) executes 4 comprehensive test cases simulating $100\text{ MHz}$ destination clock frequency ($T = 10\text{ ns}$).

\lstinputlisting[language=Verilog, caption={RST\_SYNC\_tb.v -- Self-Checking Testbench}]{../tb/RST_SYNC_tb.v}

\subsection{ModelSim Console Transcript}

All 4 test cases passed successfully with 0 errors:

\begin{figure}[H]
    \centering
    \includegraphics[width=\textwidth]{log output.png}
    \caption{ModelSim Transcript Log -- All 4 Test Cases PASSED (0 Errors)}
\end{figure}

\subsection{Output Waveform}

\begin{figure}[H]
    \centering
    \includegraphics[width=\textwidth]{output waveform.png}
    \caption{ModelSim Output Waveform -- RST\_SYNC Simulation}
\end{figure}

\subsection{Detailed Waveform and Signal Transition Analysis}

The output waveform captures the dynamic behavior of both 2-stage and 4-stage reset synchronizers:

\subsubsection{Clock and Asynchronous Reset Input}
\begin{itemize}
    \item \textbf{\texttt{CLK} (Cyan):} Destination domain clock operating at $f = 100$ MHz ($T = 10$ ns).
    \item \textbf{\texttt{RST} (Orange, Async Input):} Asynchronous active-low reset input.
    \begin{itemize}
        \item At initial time $t = 0$, \texttt{RST} $= 0$ (active reset).
        \item At $t \approx 25$ ns, \texttt{RST} is de-asserted to $1$ asynchronously mid-cycle.
        \item At $t \approx 75$ ns, \texttt{RST} is re-asserted to $0$ asynchronously mid-cycle.
        \item At $t \approx 105$ ns, \texttt{RST} is de-asserted to $1$ again mid-cycle.
    \end{itemize}
\end{itemize}

\subsubsection{2-Stage Synchronizer Operation (\texttt{NUM\_STAGES} = 2)}
\begin{itemize}
    \item \textbf{\texttt{sync\_reg} (Magenta, Binary):}
    \begin{enumerate}
        \item During active reset ($\texttt{RST} = 0$), \texttt{sync\_reg} $= \texttt{2'b00}$.
        \item When \texttt{RST} de-asserts to $1$ mid-cycle, \texttt{sync\_reg} stays \texttt{2'b00} until the first clock edge.
        \item On 1st \texttt{posedge CLK}: Stage 1 captures $1$, so \texttt{sync\_reg} $\to \texttt{2'b01}$. Output \texttt{SYNC\_RST} stays $0$.
        \item On 2nd \texttt{posedge CLK}: Stage 2 captures $1$, so \texttt{sync\_reg} $\to \texttt{2'b11}$. Output \texttt{SYNC\_RST} de-asserts to $1$ \textbf{strictly on the clock edge}.
    \end{enumerate}
    \item \textbf{\texttt{SYNC\_RST\_2stage} (Lime):} Remains $0$ during the 2-cycle synchronization latency, then rises to $1$ synchronously on the 2nd rising clock edge. When \texttt{RST} drops to $0$ asynchronously, \texttt{SYNC\_RST\_2stage} drops to $0$ \textbf{immediately} ($< 1\text{ ns}$).
\end{itemize}

\subsubsection{4-Stage Synchronizer Operation (\texttt{NUM\_STAGES} = 4)}
\begin{itemize}
    \item \textbf{\texttt{sync\_reg} (Purple, Binary):} Shifts $1\text{'b1}$ sequentially through 4 flip-flops:
    \[
        \texttt{4'b0000} \;\to\; \texttt{4'b0001} \;\to\; \texttt{4'b0011} \;\to\; \texttt{4'b0111} \;\to\; \texttt{4'b1111}
    \]
    \item \textbf{\texttt{SYNC\_RST\_4stage} (Green):} De-asserts to $1$ on the 4th rising clock edge ($40\text{ ns}$ latency), confirming that \texttt{NUM\_STAGES} parameterization correctly scales the de-assertion latency.
\end{itemize}

% ===========================================================================
% 6. SPYGLASS LINT RESULTS
% ===========================================================================
\section{SpyGlass Lint Results}

The design was analyzed using Synopsys SpyGlass (version vL-2016.06) under the \texttt{rtl\_handoff} methodology.

\subsection{Design Read Step Output}

SpyGlass successfully read and elaborated \texttt{RST\_SYNC.v}. The message tree shows \textbf{Total: 2 messages, Displayed: 2, Waived: 0} -- both are informational messages.

\begin{figure}[H]
    \centering
    \includegraphics[width=\textwidth]{Design Read Step output.png}
    \caption{SpyGlass Design Read Step Output -- 2 Info Messages, 0 Errors}
\end{figure}

\subsection{Lint RTL Goal Run Output}

The Lint RTL goal run completed with a \textbf{clean lint result} (0 Errors, 0 Warnings, 0 Latches).

\begin{figure}[H]
    \centering
    \includegraphics[width=\textwidth]{goal run output (lint clean).png}
    \caption{SpyGlass Lint RTL Goal Run Output -- Clean Lint (0 Errors, 0 Warnings)}
\end{figure}

\subsection{SpyGlass Lint Summary Report (moresimple.rpt)}

\begin{lstlisting}[language={}, caption={SpyGlass moresimple.rpt -- Lint Summary Report}]
Total Number of Generated Messages :          2
Number of Waived Messages          :          0
Number of Reported Messages        :          2
Number of Overlimit Messages       :          0

MORESIMPLE REPORT:

############### BuiltIn -> RuleGroup=Design Read ###############
[1] DetectTopDesignUnits  Info  ../rtl/RST_SYNC.v  Line 11
    Module RST_SYNC is a top level design unit
[0] ElabSummary           Info
    Please refer file elab_summary.rpt for elab summary report
\end{lstlisting}

% ===========================================================================
% 7. CONCLUSION
% ===========================================================================
\section{Conclusion}

The Reset Synchronizer (\texttt{RST\_SYNC}) module was successfully designed, modeled, and verified.
\begin{itemize}
    \item Immediate asynchronous reset assertion guarantees rapid system protection ($< 1\text{ ns}$).
    \item Synchronous reset de-assertion eliminates recovery and removal timing violations across the destination clock domain.
    \item Parameterized flip-flop stage depth (\texttt{NUM\_STAGES}) allows configurable synchronization latency.
    \item All 4 test cases passed with 0 errors in ModelSim simulation.
    \item Clean SpyGlass lint report: 0 errors, 0 warnings, 0 latches.
\end{itemize}

\end{document}
